\chapter{Química supramolecular}\label{ch:b3:supramolecular}

En 1962, Charles Pedersen, que buscaba otra cosa, aisló un poliéter cíclico
que disolvía el permanganato de potasio en benceno y lo teñía de
púrpura. El anillo se había enrollado alrededor del ion potasio y había ocultado su
carga tras una superficie hidrocarbonada, arrastrando consigo el permanganato como
contraión. La química más allá de la molécula, la química de los conjuntos mantenidos
por fuerzas más débiles que los enlaces covalentes, nació con aquella disolución púrpura.
Este capítulo mide esas fuerzas, explica por qué algunos \omterm{def:b3:supramolecular:supramolecular}{anfitriones} eligen tan bien a sus \omterm{def:b3:supramolecular:supramolecular}{huéspedes},
muestra cómo se miden las constantes de asociación y termina con moléculas
enhebradas y entrelazadas para formar máquinas.

\begin{recall}
El volumen del primer año trató las fuerzas intermoleculares, el enlace de hidrógeno, la solvatación,
las moléculas hidrófobas, los complejos, las constantes de formación y los quelatos; el volumen del
segundo año, el potencial químico, las energías de Gibbs y el ajuste por mínimos cuadrados.
El \cref{ch:b3:advanced-nmr} describió el intercambio rápido y lento y la coalescencia;
el \cref{ch:b3:rate-theories}, la \omterm{thm:b3:rate-theories:eyring}{ecuación de Eyring}, y
el \cref{ch:b3:partition-functions}, la entropía de traslación de un gas.
\end{recall}

\section{Interacciones no covalentes}

\begin{definition}[Química supramolecular]\label{def:b3:supramolecular:supramolecular}
La \emph{química supramolecular}\index{quimica supramolecular@química supramolecular} es la química
de los conjuntos de moléculas mantenidos por interacciones no covalentes. En un
\emph{complejo anfitrión--huésped}\index{complejo anfitrion--huesped@complejo anfitrión--huésped}, el \emph{anfitrión}\index{anfitrion@anfitrión}
es la molécula más grande, con una cavidad o una hendidura y sitios de unión que convergen
hacia ella, y el \emph{huésped}\index{huesped@huésped}, la molécula o el ion más pequeño que une.
\end{definition}

Las interacciones son las del volumen del primer año, ion--dipolo, dipolo--dipolo,
enlaces de hidrógeno y fuerzas de London, más algunas con nombre propio.

\begin{definition}[Interacciones no covalentes]\label{def:b3:supramolecular:interactions}
El \emph{apilamiento $\pi$}\index{apilamiento $\pi$} es la atracción entre anillos
aromáticos planos situados cara a cara, normalmente desplazados o inclinados. Una \emph{interacción
catión--$\pi$}\index{interaccion cation--@interacción catión--$\pi$} es la atracción entre un catión y
la cara rica en electrones de un anillo aromático. Un \emph{enlace de halógeno}\index{enlace de halogeno@enlace de halógeno}
es la atracción entre la punta pobre en electrones de un átomo de halógeno unido covalentemente,
opuesta a su enlace, y una base de Lewis. El \emph{efecto
hidrófobo}\index{efecto hidrofobo@efecto hidrófobo} es la tendencia de las superficies apolares a
juntarse en el agua, impulsada sobre todo por la liberación de las moléculas de agua ordenadas
que las rodean.
\end{definition}

Las interacciones ion--dipolo y catión--$\pi$ son las más fuertes, y dependen
mucho del disolvente, que compite por el ion; les siguen los enlaces de hidrógeno,
después el apilamiento $\pi$ y los \omterm{def:b3:supramolecular:interactions}{enlaces de halógeno}, y después las fuerzas de London, débiles una a una pero
sumadas sobre grandes superficies de contacto. En el agua, el \omterm{def:b3:supramolecular:interactions}{efecto hidrófobo} suele
dominar: un \omterm{def:b3:supramolecular:supramolecular}{anfitrión} diseñado para unir un \omterm{def:b3:supramolecular:supramolecular}{huésped} mediante enlaces de hidrógeno en cloroformo puede no
unirlo en absoluto en agua, cuyos propios enlaces de hidrógeno compiten.

\section{Reconocimiento molecular}

\begin{definition}[Reconocimiento molecular]\label{def:b3:supramolecular:recognition}
El \emph{reconocimiento molecular}\index{reconocimiento molecular} es la unión
selectiva de un \omterm{def:b3:supramolecular:supramolecular}{huésped} por un \omterm{def:b3:supramolecular:supramolecular}{anfitrión}. Se basa en la \emph{complementariedad}\index{complementariedad},
una correspondencia de tamaño, forma y sitios de unión entre \omterm{def:b3:supramolecular:supramolecular}{anfitrión} y \omterm{def:b3:supramolecular:supramolecular}{huésped}, y se
refuerza con la \emph{preorganización}\index{preorganizacion@preorganización}: un \omterm{def:b3:supramolecular:supramolecular}{anfitrión} cuyos
sitios de unión ya están sujetos en la geometría del complejo antes de la unión.
\end{definition}

La constante de asociación $K_a = [\mathrm{HG}]/([\mathrm H][\mathrm G])$ es una constante de formación en el sentido
del volumen del primer año, y $\Delta G^\circ = -RT\ln K_a$. Un \omterm{def:b3:supramolecular:supramolecular}{anfitrión} flexible debe congelar su
conformación al unirse, lo que cuesta entropía; uno rígido y preorganizado ya ha
pagado ese precio en su síntesis.

\begin{definition}[Efectos quelato y macrocíclico]\label{def:b3:supramolecular:macrocyclic}
El \emph{efecto quelato}\index{efecto quelato} es la mayor estabilidad de un
complejo con un ligando polidentado que con el número equivalente de
ligandos monodentados con los mismos átomos dadores. El \emph{efecto
macrocíclico}\index{efecto macrociclico@efecto macrocíclico} es la estabilización adicional cuando el
ligando polidentado es un anillo y no una cadena abierta.
\end{definition}

\begin{proposition}[El efecto quelato es sobre todo entrópico]\label{prop:b3:supramolecular:chelate-entropy}
En $\mathrm{[M(NH_3)_6]^{2+} + 3\,en \rightleftharpoons [M(en)_3]^{2+} + 6\,NH_3}$, donde en es la etano-1,2-diamina, los mismos seis
enlaces metal--nitrógeno están presentes a ambos lados, y la reacción la impulsa
el aumento del número de moléculas libres.
\end{proposition}

\begin{proof}[Argumento]
Los enlaces formados y rotos son del mismo tipo, de modo que $\Delta_rH^\circ$ es pequeña. El número de
moléculas independientes pasa de cuatro a siete: tres grados de libertad de traslación
más de una molécula entera, que valen cada uno decenas de julios por kelvin
y por mol en disolución (\cref{ch:b3:partition-functions}), luego $\Delta_rS^\circ > 0$ y $K > 1$.
Desde el punto de vista cinético, una vez unido un extremo de un quelato, el otro se mantiene cerca del
metal a una concentración efectiva alta, y vuelve a unirse enseguida si se suelta.
\end{proof}

\begin{proposition}[Compensación entalpía--entropía]\label{prop:b3:supramolecular:compensation}
En una serie de \omterm{def:b3:supramolecular:supramolecular}{complejos anfitrión--huésped} relacionados, un $\Delta H^\circ$ de unión más negativo
suele ir acompañado de un $\Delta S^\circ$ más negativo, de modo que $\Delta G^\circ$ varía menos que cualquiera de los dos (una
observación experimental).
\end{proposition}

Un complejo más apretado se une con más fuerza, pero restringe más los movimientos del \omterm{def:b3:supramolecular:supramolecular}{anfitrión} y
del \omterm{def:b3:supramolecular:supramolecular}{huésped}; un enlace de hidrógeno más fuerte con el \omterm{def:b3:supramolecular:supramolecular}{huésped} suele significar uno más débil con
el agua que desplazó. La energía de Gibbs, que fija la selectividad, es la pequeña
diferencia de dos términos grandes.

\begin{definition}[Anfitriones]\label{def:b3:supramolecular:hosts}
Un \emph{éter corona}\index{eter corona@éter corona} es un poliéter macrocíclico, como el
18-corona-6, $\ce{(CH2CH2O)6}$. Un \emph{criptando}\index{criptando} es un poliéter
bicíclico con dos cabezas de puente de nitrógeno, que encierra un catión en tres
dimensiones en un \emph{criptato}\index{criptato}. Una
\emph{ciclodextrina}\index{ciclodextrina} es un oligómero cíclico de unidades de glucosa, con
forma de cono truncado, una cavidad hidrófoba y grupos hidroxi en sus
bordes. Un \emph{calixareno}\index{calixareno} es un macrociclo en forma de copa formado por
unidades de fenol unidas por puentes metileno.
\end{definition}

\begin{omfigure}
\begin{tikzpicture}[font=\small]
  % 18-crown-6 with K+
  \foreach \k in {0,...,17} {
    \pgfmathsetmacro\ang{90 + 20*\k}
    \pgfmathsetmacro\nxt{110 + 20*\k}
    \draw[black!70] (\ang:1.55) -- (\nxt:1.55);
  }
  \foreach \k in {0,3,...,15} {
    \pgfmathsetmacro\ang{90 + 20*\k}
    \node[fill=white, inner sep=0.5pt, text=omThm] (o\k) at (\ang:1.55) {O};
    \draw[omThm!60, densely dashed] (0,0) -- (\ang:1.32);
  }
  \node[circle, fill=omDef!25, draw=omDef, inner sep=1pt] at (0,0) {$\ce{K+}$};
  \node[font=\scriptsize] at (0,-2.1) {18-corona-6 con $\ce{K+}$};
  % [2.2.2]cryptand with K+
  \begin{scope}[xshift=5.2cm]
    \node (nt) at (0,1.7) {N};
    \node (nb) at (0,-1.7) {N};
    \foreach \x/\n in {-1.25/a, 0/b, 1.25/c} {
      \node[text=omThm, fill=white, inner sep=0.5pt] (oa\n) at (\x,0.55) {O};
      \node[text=omThm, fill=white, inner sep=0.5pt] (ob\n) at (\x,-0.55) {O};
      \draw[black!70] (nt) to[out=-90+40*\x, in=90] (oa\n) -- (ob\n) to[out=-90, in=90-40*\x] (nb);
    }
    \node[circle, fill=omDef!25, draw=omDef, inner sep=1pt] at (0.62,0) {$\ce{K+}$};
    \node[font=\scriptsize] at (0,-2.4) {[2.2.2]criptando: criptato de $\ce{K+}$};
  \end{scope}
\end{tikzpicture}

{\small Izquierda: el 18-corona-6, un anillo de doce átomos de carbono (vértices) y seis átomos de
oxígeno, con un ion potasio en su centro unido por los seis oxígenos. Derecha: el
[2.2.2]\omterm{def:b3:supramolecular:hosts}{criptando}, tres cadenas de dos oxígenos de éter entre dos cabezas de puente
de nitrógeno, en esquema: el catión queda sujeto dentro de una jaula tridimensional,
por seis oxígenos y los dos nitrógenos.}
\end{omfigure}

Los \omterm{def:b3:supramolecular:hosts}{éteres corona} seleccionan los cationes por su tamaño: el 18-corona-6 une el $\ce{K+}$ con más fuerza que el
$\ce{Na+}$, más pequeño, que no puede alcanzar a la vez los seis oxígenos, o que el $\ce{Cs+}$, más grande, que
se sitúa por encima del anillo. Los \omterm{def:b3:supramolecular:hosts}{criptandos}, más preorganizados y que envuelven el ion
por completo, se unen con más fuerza aún y seleccionan con más nitidez. Oculto dentro
del \omterm{def:b3:supramolecular:supramolecular}{anfitrión}, el catión arrastra a su anión hacia los disolventes apolares, donde el anión,
no solvatado, se vuelve muy reactivo: la base de la catálisis por transferencia de fase.

\section{Medir la unión}

\begin{proposition}[Intercambio rápido]\label{prop:b3:supramolecular:fast-exchange}
Cuando un \omterm{def:b3:supramolecular:supramolecular}{anfitrión} se intercambia entre sus estados libre y unido mucho más deprisa que la
diferencia de sus frecuencias de resonancia, su señal de RMN aparece en la media
$\delta_{\mathrm{obs}} = x_{\mathrm{libre}}\delta_{\mathrm{libre}} + x_{\mathrm{unido}}\delta_{\mathrm{unido}}$, donde las $x$ son las fracciones de \omterm{def:b3:supramolecular:supramolecular}{anfitrión} en cada estado.
\end{proposition}

\begin{proof}
Cada molécula de \omterm{def:b3:supramolecular:supramolecular}{anfitrión} visita ambos estados muchas veces durante la medida; su
frecuencia de precesión, promediada en el tiempo, es la media ponderada de las dos, con
pesos iguales a las fracciones de tiempo pasadas en cada uno, que en el equilibrio son las
fracciones de moléculas en cada estado (\cref{ch:b3:advanced-nmr}).
\end{proof}

\begin{theorem}[Isoterma 1:1 exacta]\label{thm:b3:supramolecular:isotherm}
Para $\mathrm{H + G \rightleftharpoons HG}$, con concentraciones totales $H_0$ y $G_0$ y constante de asociación $K$,
\[
  [\mathrm{HG}] = \frac{1}{2}\Bigl[\Bigl(H_0 + G_0 + \frac1K\Bigr) - \sqrt{\Bigl(H_0 + G_0 + \frac1K\Bigr)^2 - 4H_0G_0}\,\Bigr].
\]
\end{theorem}

\begin{proof}
Con $c = [\mathrm{HG}]$, $K = c/((H_0 - c)(G_0 - c))$, es decir, $c^2 - (H_0 + G_0 + 1/K)c + H_0G_0 = 0$. De las dos raíces,
solo la menor es inferior a la vez a $H_0$ y a $G_0$ (su producto es $H_0G_0$ y su suma
supera $H_0 + G_0$): es la del signo menos.
\end{proof}

\begin{omfigure}
\begin{tikzpicture}
\begin{axis}[name=a, width=0.48\linewidth, height=5.5cm, xmin=0, xmax=6, ymin=0, ymax=1.05,
  axis lines=left, xlabel={equivalentes de huésped, $G_0/H_0$}, ylabel={fracción de anfitrión unido},
  tick label style={font=\scriptsize}, label style={font=\small},
  ]
\addplot[omMeth, thick] table[x=equiv, y=K4] {figdata/out/bachelor-3-binding-titration.dat};
\addplot[omDef, thick] table[x=equiv, y=K3] {figdata/out/bachelor-3-binding-titration.dat};
\addplot[omThm, thick] table[x=equiv, y=K2] {figdata/out/bachelor-3-binding-titration.dat};
\node[font=\scriptsize, omMeth, anchor=north] at (axis cs:4,0.94) {$K = 10^4$};
\node[font=\scriptsize, omDef, anchor=north] at (axis cs:4,0.68) {$K = 10^3$};
\node[font=\scriptsize, omThm, anchor=north] at (axis cs:4,0.27) {$K = 10^2$};
\end{axis}
\begin{axis}[at={(a.east)}, anchor=west, xshift=1.4cm, width=0.48\linewidth, height=5.5cm,
  xmin=0, xmax=1, ymin=0, ymax=1.05, axis lines=left,
  xlabel={fracción molar de anfitrión}, ylabel={$[\text{complejo}]/c_{\text{total}}$},
  tick label style={font=\scriptsize}, label style={font=\small}]
\addplot[omDef, thick] table[x=x, y=HG] {figdata/out/bachelor-3-binding-titration-b.dat};
\addplot[omThm, thick] table[x=x, y=HG2] {figdata/out/bachelor-3-binding-titration-b.dat};
\draw[black!35, dashed] (axis cs:0.5,0) -- (axis cs:0.5,0.55) (axis cs:0.3333,0) -- (axis cs:0.3333,0.38);
\node[font=\scriptsize, omDef, anchor=south] at (axis cs:0.5,0.52) {HG};
\node[font=\scriptsize, omThm, anchor=south] at (axis cs:0.25,0.36) {$\mathrm{HG_2}$};
\end{axis}
\end{tikzpicture}

{\small Izquierda: fracción de un \omterm{def:b3:supramolecular:supramolecular}{anfitrión} unido durante una valoración con $H_0 = \qty{1.0}{mmol/L}$ para tres
constantes de asociación (en $\unit{L.mol^{-1}}$): cuanto mayor es $K$, más neto es el quiebro en un equivalente;
cuando $KH_0 \gg 1$, la curva solo indica que $K$ es grande. Derecha: \omterm{def:b3:supramolecular:job}{representaciones de Job}, la
concentración de complejo a concentración total constante de \omterm{def:b3:supramolecular:supramolecular}{anfitrión} más \omterm{def:b3:supramolecular:supramolecular}{huésped},
con un máximo en una fracción de \omterm{def:b3:supramolecular:supramolecular}{anfitrión} de 1/2 para un complejo 1:1 y de 1/3 para uno 1:2 (constantes
del modelo).}
\end{omfigure}

\begin{method}[Constante de asociación a partir de una valoración por RMN]\label{met:b3:supramolecular:nmr-titration}
\begin{enumerate}
  \item Mantener fija la concentración de \omterm{def:b3:supramolecular:supramolecular}{anfitrión} $H_0$ y añadir \omterm{def:b3:supramolecular:supramolecular}{huésped}, de 0 a varios
        equivalentes; registrar una señal del \omterm{def:b3:supramolecular:supramolecular}{anfitrión} en cada punto (intercambio rápido).
  \item Modelo: $\delta = \delta_{\mathrm{libre}} + (\delta_{\mathrm{unido}} - \delta_{\mathrm{libre}})[\mathrm{HG}]/H_0$, con $[\mathrm{HG}]$ dada por la isoterma exacta.
  \item Ajustar $K$ y $\delta_{\mathrm{unido}}$ por mínimos cuadrados no lineales, minimizando la suma de los cuadrados
        de los residuos; tomar los errores típicos de la curvatura de esa suma
        (la matriz de covarianza).
  \item Comprobar que los residuos no muestran tendencia (estequiometría errónea) y que la
        fracción unida abarca aproximadamente del 20 al 80\,\% (los datos acotan entonces $K$).
\end{enumerate}
\end{method}

\begin{definition}[Representación de Job]\label{def:b3:supramolecular:job}
Una \emph{representación de Job}\index{representacion de Job@representación de Job} (método de las variaciones continuas) es una representación de
la concentración de un complejo, o de una señal proporcional a ella, en función de
la fracción molar de uno de los componentes, a concentración total constante de los
dos.
\end{definition}

\begin{proposition}[Máximo de una representación de Job]\label{prop:b3:supramolecular:job}
Para un único complejo $\mathrm{H_nG_m}$, la \omterm{def:b3:supramolecular:job}{representación de Job} es máxima en la fracción molar de \omterm{def:b3:supramolecular:supramolecular}{anfitrión}
$x = n/(n + m)$, sea cual sea la constante de asociación.
\end{proposition}

\begin{proof}
Con concentración total $T$, $H_t = xT$, $G_t = (1 - x)T$, complejo $c$, libres $[\mathrm H] = H_t - nc$ y
$[\mathrm G] = G_t - mc$, y $c = \beta[\mathrm H]^n[\mathrm G]^m$. Derivando respecto de $x$ en el máximo, donde
$\dd c/\dd x = 0$: $0 = \beta\bigl(n[\mathrm H]^{n-1}[\mathrm G]^m T - m[\mathrm H]^n[\mathrm G]^{m-1}T\bigr)$, luego $n[\mathrm G] = m[\mathrm H]$. Sustituyendo,
$n(G_t - mc) = m(H_t - nc)$, luego $nG_t = mH_t$, es decir, $n(1 - x) = mx$: $x = n/(n + m)$.
\end{proof}

\begin{method}[Realizar una representación de Job]\label{met:b3:supramolecular:job}
\begin{enumerate}
  \item Preparar disoluciones madre equimolares de \omterm{def:b3:supramolecular:supramolecular}{anfitrión} y de \omterm{def:b3:supramolecular:supramolecular}{huésped}.
  \item Mezclarlas en proporciones $x$ de 0 a 1 con un volumen total constante.
  \item Medir una propiedad proporcional al complejo (una absorbancia corregida
        de las especies libres, o $x\,\Delta\delta$ en RMN).
  \item Leer la estequiometría en la posición del máximo.
\end{enumerate}
\end{method}

\begin{definition}[Calorimetría de valoración isotérmica]\label{def:b3:supramolecular:itc}
La \emph{calorimetría de valoración isotérmica}\index{calorimetria de valoracion isotermica@calorimetría de valoración isotérmica}
(ITC) mide el calor liberado o absorbido en cada inyección de una disolución de \omterm{def:b3:supramolecular:supramolecular}{huésped}
en una disolución de \omterm{def:b3:supramolecular:supramolecular}{anfitrión} mantenida a temperatura constante.
\end{definition}

\begin{proposition}[Calor por inyección]\label{prop:b3:supramolecular:itc}
Despreciando la dilución, el calor de la $i$-ésima inyección en una celda de volumen $V_0$ es
$q_i = \Delta H^\circ V_0\bigl([\mathrm{HG}]_i - [\mathrm{HG}]_{i-1}\bigr)$.
\end{proposition}

\begin{proof}
El calor es la entalpía de reacción multiplicada por la cantidad de complejo formado durante la
inyección, y esa cantidad es la variación de $[\mathrm{HG}]$ multiplicada por el volumen de la celda.
\end{proof}

\begin{method}[Leer una isoterma de ITC]\label{met:b3:supramolecular:itc}
\begin{enumerate}
  \item Representar el calor por mol de \omterm{def:b3:supramolecular:supramolecular}{huésped} inyectado en función de la razón molar.
  \item La altura de la meseta antes de la saturación da $\Delta H^\circ$; la razón molar en
        el punto de inflexión da la estequiometría; la pendiente da $K$.
  \item Ajustar los tres con la isoterma exacta; después, $\Delta G^\circ = -RT\ln K$ y
        $T\Delta S^\circ = \Delta H^\circ - \Delta G^\circ$.
\end{enumerate}
\end{method}

Un solo experimento de ITC da así toda la firma termodinámica de la
unión, entalpía y entropía, que una valoración por RMN o UV--visible solo da
a través de una serie de temperaturas.

\begin{inthelab}[Una valoración por RMN]
Una disolución del \omterm{def:b3:supramolecular:supramolecular}{anfitrión} en un disolvente deuterado, de un milimol por litro aproximadamente,
se coloca en un tubo de RMN y se registra su espectro. Una disolución concentrada del
\omterm{def:b3:supramolecular:supramolecular}{huésped} preparada en la disolución del \omterm{def:b3:supramolecular:supramolecular}{anfitrión}, de modo que la concentración de \omterm{def:b3:supramolecular:supramolecular}{anfitrión} permanezca
constante, se añade en pequeñas porciones con una microjeringa; tras cada adición, el
tubo se agita, se equilibra a la temperatura de la sonda y se registra el
espectro. El desplazamiento de una señal del \omterm{def:b3:supramolecular:supramolecular}{anfitrión} que se mueve se lee en cada punto y se
ajusta.
\end{inthelab}

\section{Anfitriones}

Las \omterm{def:b3:supramolecular:hosts}{ciclodextrinas}, obtenidas por vía enzimática a partir del almidón, tienen seis, siete u ocho unidades
de glucosa ($\alpha$, $\beta$, $\gamma$). Su cavidad hidrófoba aloja las partes apolares de los \omterm{def:b3:supramolecular:supramolecular}{huéspedes} en
agua: solubilizan fármacos, transportan fragancias y extraen el colesterol de las
membranas celulares. Los \omterm{def:b3:supramolecular:hosts}{calixarenos} y los cucurbiturilos, con forma de calabaza, macrociclos
rígidos de unidades de glicolurilo con portales revestidos de carbonilos, unen cationes y
cationes orgánicos muy fuertemente en agua.

\begin{omfigure}
\begin{tikzpicture}[font=\small]
  \draw[black!70, fill=omDef!8] (-1.9,1.2) -- (-1.2,-1.2) arc[start angle=180, end angle=360, x radius=1.2, y radius=0.3]
    -- (1.9,1.2);
  \draw[black!70, fill=omDef!18] (0,1.2) ellipse[x radius=1.9, y radius=0.45];
  \draw[black!40, dashed] (-1.2,-1.2) arc[start angle=180, end angle=0, x radius=1.2, y radius=0.3];
  \node[draw=omThm, thick, regular polygon, regular polygon sides=6, minimum size=8mm, inner sep=0pt] at (0,0.2) {};
  \draw[omThm, thick] (0,0.62) -- (0,2.0);
  \node[font=\scriptsize, omThm, anchor=west] at (0.1,2.0) {huésped};
  \node[font=\scriptsize, anchor=west] at (2.0,1.2) {borde ancho: OH secundarios};
  \node[font=\scriptsize, anchor=west] at (1.3,-1.25) {borde estrecho: OH primarios};
  \node[font=\scriptsize, anchor=east] at (-1.75,0) {cavidad hidrófoba};
\end{tikzpicture}

{\small Una \omterm{def:b3:supramolecular:hosts}{ciclodextrina}, dibujada como el cono truncado que forma, con un
\omterm{def:b3:supramolecular:supramolecular}{huésped} aromático en su cavidad. Los grupos hidroxi de los dos bordes hacen el exterior
soluble en agua; el interior, revestido de grupos C--H y de oxígenos glicosídicos, es
apolar.}
\end{omfigure}

\section{Autoensamblaje y máquinas moleculares}

\begin{definition}[Autoensamblaje]\label{def:b3:supramolecular:self-assembly}
El \emph{autoensamblaje}\index{autoensamblaje} es la organización espontánea y reversible
de componentes en una estructura definida, bajo el control de la
información contenida en los propios componentes: sus formas y sus sitios
de unión.
\end{definition}

Los iones metálicos con ángulos de coordinación fijos y los ligandos puente rígidos se ensamblan
en cuadrados, jaulas y cápsulas en una sola etapa, porque cada enlace puede abrirse
y cerrarse hasta que se alcanza la estructura más estable, sin tensión ni posiciones
sueltas. El apareamiento de bases en el ADN y el plegamiento de las proteínas son la misma idea
a mayor escala.

\begin{definition}[Moléculas mecánicamente entrelazadas]\label{def:b3:supramolecular:interlocked}
Una \emph{molécula mecánicamente entrelazada}\index{molecula mecanicamente entrelazada@molécula mecánicamente entrelazada}
tiene partes que no están enlazadas entre sí pero que no pueden separarse
sin romper un enlace covalente. Un \emph{rotaxano}\index{rotaxano} es un anillo
enhebrado en un eje con topes voluminosos en ambos extremos; un
\emph{catenano}\index{catenano} consta de anillos entrelazados.
\end{definition}

\begin{definition}[Síntesis con plantilla]\label{def:b3:supramolecular:template}
Una \emph{síntesis con plantilla}\index{sintesis con plantilla@síntesis con plantilla} usa una interacción
temporal, a menudo con un ion metálico, para mantener los componentes reactivos en la
geometría necesaria para una reacción que de otro modo sería improbable, como el
cierre de un anillo alrededor de otro anillo.
\end{definition}

\begin{omfigure}
\begin{tikzpicture}[font=\small, >=Stealth]
  % step 1: two crescents around Cu+
  \draw[omDef, line width=2pt] (0.3,0) arc[start angle=0, end angle=300, radius=0.75];
  \draw[omThm, line width=2pt] (-0.3,0) arc[start angle=180, end angle=-120, radius=0.75];
  \fill[omProp] (0,0) circle (0.13);
  \node[font=\scriptsize] at (0,-1.35) {el $\ce{Cu+}$ sujeta dos};
  \node[font=\scriptsize] at (0,-1.7) {ligandos abiertos cruzados};
  \draw[->, thick] (1.4,0) -- (2.3,0) node[midway, above, font=\scriptsize] {cierre};
  % step 2: closed interlocked rings with Cu
  \begin{scope}[xshift=3.6cm]
    \draw[omDef, line width=2pt] (-0.35,0) circle (0.75);
    \draw[omThm, line width=2pt] (0.35,0) circle (0.75);
    \draw[white, line width=5pt] ([shift={(52:0.75)}]-0.35,0) arc[start angle=52, end angle=72, radius=0.75];
    \draw[omDef, line width=2pt] ([shift={(48:0.75)}]-0.35,0) arc[start angle=48, end angle=76, radius=0.75];
    \fill[omProp] (0,0) circle (0.13);
    \node[font=\scriptsize] at (0,-1.35) {catenato};
  \end{scope}
  \draw[->, thick] (5.0,0) -- (5.9,0) node[midway, above, font=\scriptsize] {$-\ce{Cu+}$};
  \begin{scope}[xshift=7.2cm]
    \draw[omDef, line width=2pt] (-0.35,0) circle (0.75);
    \draw[omThm, line width=2pt] (0.35,0) circle (0.75);
    \draw[white, line width=5pt] ([shift={(52:0.75)}]-0.35,0) arc[start angle=52, end angle=72, radius=0.75];
    \draw[omDef, line width=2pt] ([shift={(48:0.75)}]-0.35,0) arc[start angle=48, end angle=76, radius=0.75];
    \node[font=\scriptsize] at (0,-1.35) {[2]catenano};
  \end{scope}
  % rotaxane
  \begin{scope}[xshift=10.2cm]
    \draw[black!70, line width=1.6pt] (-1.1,0) -- (1.1,0);
    \fill[black!60] (-1.25,0) circle (0.2) (1.25,0) circle (0.2);
    \draw[omThm, line width=2pt] (0,0) ellipse[x radius=0.18, y radius=0.55];
    \node[font=\scriptsize] at (0,-1.35) {[2]rotaxano};
  \end{scope}
\end{tikzpicture}

{\small \omterm{def:b3:supramolecular:template}{Síntesis con plantilla} de un \omterm{def:b3:supramolecular:interlocked}{catenano} (esquema): un ion cobre(I) sujeta dos
ligandos derivados de la fenantrolina cruzados en ángulo recto en su coordinación
tetraédrica; cerrar cada uno en un anillo da dos anillos entrelazados alrededor del
metal, y eliminar el cobre deja el \omterm{def:b3:supramolecular:interlocked}{catenano} libre. Derecha: un \omterm{def:b3:supramolecular:interlocked}{rotaxano}, un
anillo atrapado en un eje por dos topes.}
\end{omfigure}

\begin{definition}[Máquinas moleculares]\label{def:b3:supramolecular:machine}
Una \emph{máquina molecular}\index{maquina molecular@máquina molecular} es un conjunto cuyos
componentes se mueven unos respecto de otros de forma controlada en respuesta a un
estímulo (la luz, un cambio redox, un cambio de pH), y al que se puede hacer realizar un trabajo.
\end{definition}

En una lanzadera molecular, un anillo de un eje de \omterm{def:b3:supramolecular:interlocked}{rotaxano} se sitúa en una de dos estaciones; un
estímulo que debilita su unión a la primera lo envía a la segunda, y el
estímulo inverso lo devuelve. Los motores rotatorios impulsados por la luz basados en
alquenos sobrecargados estéricamente giran en un solo sentido, alternando \omterm{def:b3:photochemistry:photoisomerisation}{fotoisomerizaciones}
y etapas térmicas que no pueden ir hacia atrás. La velocidad de
desplazamiento entre estaciones se mide por RMN a temperatura variable, a partir de la
coalescencia de las señales de las dos formas (\cref{ch:b3:advanced-nmr}).

\begin{safety}{\ghs[0.62cm]{GHS07}}
% ledger: ghs4:18C6
18-Corona-6: nocivo en caso de ingestión, irritante para la piel, los ojos y las vías respiratorias.
Los \omterm{def:b3:supramolecular:hosts}{éteres corona} transportan sales metálicas hacia los disolventes orgánicos y a través de la piel; se
pesan y se manipulan con guantes.
\end{safety}

\begin{history}[La química supramolecular y sus máquinas]
Los \omterm{def:b3:supramolecular:hosts}{éteres corona} de Charles Pedersen (1967), los \omterm{def:b3:supramolecular:hosts}{criptandos} de Jean-Marie Lehn y los \omterm{def:b3:supramolecular:supramolecular}{anfitriones}
preorganizados de Donald Cram fundaron el campo; los tres compartieron el Premio Nobel
de Química de 1987. Jean-Pierre Sauvage obtuvo \omterm{def:b3:supramolecular:interlocked}{catenanos} con alto rendimiento mediante la
plantilla de cobre en 1983, Fraser Stoddart construyó lanzaderas de \omterm{def:b3:supramolecular:interlocked}{rotaxano}, y Ben
Feringa, un motor rotatorio impulsado por la luz en 1999; compartieron el Premio Nobel de
Química de 2016 por el diseño y la síntesis de \omterm{def:b3:supramolecular:machine}{máquinas moleculares}.
\end{history}

\section{Ejercicios}

\begin{exercise}[$\star$]\label{exo:b3:supramolecular:1}
Nómbrese la interacción principal en: el $\ce{K+}$ en el 18-corona-6; dos anillos de benceno cara a cara;
el $\ce{K+}$ sobre un anillo de benceno; el yodopentafluorobenceno con la piridina; un par
guanina--citosina.
\end{exercise}

\begin{exercise}[$\star$]\label{exo:b3:supramolecular:2}
Calcúlese $\Delta G^\circ$ a \qty{298}{K} para un \omterm{def:b3:supramolecular:supramolecular}{complejo anfitrión--huésped} con $K_a = \qty{1.0e4}{L.mol^{-1}}$.
\end{exercise}

\begin{exercise}[$\star$]\label{exo:b3:supramolecular:3}
En metanol, $\log K = 6.1$ para el $\ce{K+}$ y 4.4 para el $\ce{Na+}$ con el 18-corona-6 (datos del ejercicio).
Calcúlense la selectividad y la diferencia de $\Delta G^\circ$.
\end{exercise}

\begin{exercise}[$\star$]\label{exo:b3:supramolecular:4}
Una \omterm{def:b3:supramolecular:job}{representación de Job} tiene su máximo en una fracción molar de \omterm{def:b3:supramolecular:supramolecular}{anfitrión} de 0.33. ¿Cuál es la estequiometría?
\end{exercise}

\begin{exercise}[$\star\star$]\label{exo:b3:supramolecular:5}
Un \omterm{def:b3:supramolecular:supramolecular}{anfitrión} a \qty{2.0}{mmol/L} con $K = \qty{1500}{L.mol^{-1}}$, $\delta_{\mathrm{libre}} = \qty{3.560}{ppm}$ y $\delta_{\mathrm{unido}} = \qty{3.720}{ppm}$ recibe un
equivalente de \omterm{def:b3:supramolecular:supramolecular}{huésped}. Calcúlense la fracción unida y el desplazamiento observado.
\end{exercise}

\begin{exercise}[$\star\star$]\label{exo:b3:supramolecular:6}
Para el níquel(II), $\log\beta_6 = 8.6$ con el amoníaco y $\log\beta_3 = 18.3$ con la etano-1,2-diamina (datos del
ejercicio). Calcúlense la constante y el $\Delta G^\circ$ a \qty{298}{K} de
$\ce{[Ni(NH3)6]^2+} + 3\,\mathrm{en} \rightleftharpoons \ce{[Ni(en)3]^2+} + \ce{6NH3}$, y dígase qué la impulsa.
\end{exercise}

\begin{exercise}[$\star\star$]\label{exo:b3:supramolecular:7}
Un experimento de ITC a \qty{298}{K} da $K = \qty{2.0e5}{L.mol^{-1}}$ y $\Delta H^\circ = \qty{-40}{kJ/mol}$. Calcúlense $\Delta G^\circ$, $T\Delta S^\circ$ y
$\Delta S^\circ$.
\end{exercise}

\begin{exercise}[$\star\star$]\label{exo:b3:supramolecular:8}
Explíquese el papel del cobre(I), y de su geometría tetraédrica, en la síntesis de \omterm{def:b3:supramolecular:interlocked}{catenanos} de
Sauvage. ¿Cómo se elimina el cobre al final?
\end{exercise}

\begin{exercise}[$\star\star$]\label{exo:b3:supramolecular:9}
Dos señales de una lanzadera de \omterm{def:b3:supramolecular:interlocked}{rotaxano} separadas \qty{120}{Hz} coalescen a \qty{300}{K} (datos del
ejercicio). Calcúlense la constante de intercambio y la \omterm{def:b3:rate-theories:activation}{energía de Gibbs de activación}.
\end{exercise}

\begin{exercise}[$\star\star\star$]\label{exo:b3:supramolecular:10}
Un fármaco de solubilidad intrínseca \qty{0.10}{mmol/L} forma un complejo 1:1 con una
\omterm{def:b3:supramolecular:hosts}{ciclodextrina}, $K = \qty{500}{L.mol^{-1}}$ (datos del ejercicio). Demuéstrese que su solubilidad total con
\omterm{def:b3:supramolecular:hosts}{ciclodextrina} total $C_t$ es $S_0 + KS_0C_t/(1 + KS_0)$; calcúlese para $C_t = \qty{10}{mmol/L}$.
\end{exercise}

\begin{exercise}[$\star\star\star$]\label{exo:b3:supramolecular:11}
Demuéstrese, a partir de la isoterma exacta, que cuando $KH_0 \gg 1$ la fracción unida es igual al
número de equivalentes hasta uno, y después se queda en uno. ¿Por qué una valoración así
solo da una cota inferior de $K$?
\end{exercise}

\begin{exercise}[$\star\star\star$]\label{exo:b3:supramolecular:12}
Para el \omterm{def:b3:supramolecular:supramolecular}{anfitrión} del \cref{exo:b3:supramolecular:5}, de asociación
\omterm{def:b3:rate-theories:diffusion}{controlada por la difusión}, $k_{\mathrm{on}} = \qty{1e9}{L.mol^{-1}.s^{-1}}$, calcúlese $k_{\mathrm{off}}$ y demuéstrese que el intercambio es rápido
en RMN a \qty{400}{MHz}.
\end{exercise}

\section{Problema: benceno púrpura}

\begin{problem}[{Problema de fin de semana --- benceno púrpura: el potasio en un éter corona,
el permanganato extraído al tolueno, una constante de asociación medida por RMN
y por qué basta una cantidad catalítica de corona}]
\label{pb:b3:supramolecular:1}
% ledger: const:R, ghs:KMnO4, ghs:toluene, rion:K+VI, rion:Na+VI
Datos del problema. En metanol a \qty{298}{K}, $\log K = 6.1$ para el $\ce{K+}$ y 4.4 para el $\ce{Na+}$ con el
18-corona-6. Extracción: una fase acuosa con \qty{0.10}{mol/L} de $\ce{KCl}$ y \qty{1.0}{mmol/L} de $\ce{KMnO4}$ se agita
con igual volumen de tolueno con la corona L; la única especie
extraída es el par $\ce{[KL]+}\ce{MnO4-}$, con
$K_{\mathrm{ex}} = [\ce{[KL]+MnO4-}]_{\mathrm{org}}/([\ce{K+}]_{\mathrm{aq}}[\ce{MnO4-}]_{\mathrm{aq}}[\mathrm L]_{\mathrm{org}}) = \qty{1.0e5}{L^2.mol^{-2}}$, y la corona queda en el tolueno. Valoración por RMN de un
\omterm{def:b3:supramolecular:hosts}{éter corona} (\qty{2.0}{mmol/L}) con sal de sodio en metanol deuterado: señal $\ce{CH2}$
(ppm) frente a equivalentes de $\ce{Na+}$:

\begin{center}\small
\begin{tabular}{l*{11}{c}}
\hline
equiv. & 0 & 0.25 & 0.5 & 0.75 & 1.0 & 1.5 & 2.0 & 3.0 & 4.0 & 6.0 & 8.0\\
$\delta$ & 3.560 & 3.590 & 3.612 & 3.635 & 3.649 & 3.674 & 3.688 & 3.697 & 3.704 & 3.708 & 3.714\\
\hline
\end{tabular}
\end{center}
% table: binding-titration-c

\medskip\noindent\textbf{Parte I --- El complejo.}
\begin{enumerate}
  \item Dibújense el 18-corona-6 y su complejo de potasio.
  \item Calcúlese el $\Delta G^\circ$ de unión del $\ce{K+}$ y del $\ce{Na+}$.
  \item Calcúlese la selectividad $\ce{K+}/\ce{Na+}$.
  \item Con los radios en coordinación seis $r(\ce{K+}) = \qty{138}{pm}$ y $r(\ce{Na+}) = \qty{102}{pm}$, explíquese la selectividad.
  \item ¿Por qué el complejo es soluble en tolueno y el $\ce{KCl}$ no?
  \item ¿Por qué la corona se une con menos fuerza en agua que en metanol?
\end{enumerate}

\medskip\noindent\textbf{Parte II --- Extracción.}
\begin{enumerate}[resume]
  \item Escríbase el equilibrio de extracción.
  \item Siendo $y$ la concentración extraída, escríbase la ecuación de $y$ cuando
        $[\mathrm L]_0 = \qty{1.0}{mmol/L}$.
  \item Resuélvase y dese la fracción de permanganato extraída.
  \item Calcúlese la fracción con $[\mathrm L]_0 = \qty{10}{mmol/L}$.
  \item Calcúlese con $[\mathrm L]_0 = \qty{0.10}{mmol/L}$.
  \item ¿Por qué el tolueno es púrpura, y por qué el permanganato extraído es un oxidante
        más fuerte que en agua?
\end{enumerate}

\medskip\noindent\textbf{Parte III --- Una valoración por RMN.}
\begin{enumerate}[resume]
  \item ¿Por qué se desplaza una señal promediada en lugar de verse dos señales?
  \item Escríbase el modelo de $\delta$ en función de la concentración de \omterm{def:b3:supramolecular:supramolecular}{huésped}.
  \item Ajústense $K$ y $\delta_{\mathrm{unido}}$ por mínimos cuadrados; dese $K$ con su error típico.
  \item Dese la dispersión de los residuos y coméntese.
  \item ¿En qué intervalo de equivalentes está la fracción unida entre el 20 y el 80\,\%?
  \item ¿Por qué la misma valoración no permitiría medir $K$ para el potasio, $\log K = 6.1$?
\end{enumerate}

\medskip\noindent\textbf{Parte IV --- Catálisis por transferencia de fase.}
\begin{enumerate}[resume]
  \item Un alqueno en tolueno se oxida con el permanganato extraído. ¿Dónde tiene lugar
        la reacción?
  \item ¿Qué le ocurre a la corona tras la reacción, y por qué basta una cantidad
        catalítica?
  \item ¿Por qué es útil un exceso de $\ce{KCl}$ en el agua?
  \item ¿Qué catalizadores más baratos hacen el mismo trabajo en la industria?
  \item ¿Por qué se usa hoy tolueno donde Pedersen usaba benceno?
  \item Enúnciese el resultado: la fracción de permanganato extraída con
        $[\mathrm L]_0 = \qty{1.0}{mmol/L}$.
\end{enumerate}
\end{problem}
